\section{Testes e validação}\label{sec:testes-validacao}

Os testes do sistema foram planejados com o objetivo de verificar o funcionamento das principais funcionalidades implementadas, a integração entre as camadas da aplicação e a adequação do sistema às necessidades relacionadas ao apoio à avaliação psicopedagógica. A estratégia de validação foi estruturada em duas frentes: testes técnicos realizados durante o desenvolvimento e avaliação de uso com profissionais da área de psicopedagogia.

No backend, os testes automatizados (Seção~\ref{sec:backend}) cobriram dois dos níveis descritos na Seção~\ref{sec:fund_testes}: testes unitários das regras de negócio isoladas e testes de integração que exercitam as rotas da API de ponta a ponta, com o acesso ao banco de dados substituído por um objeto simulado. Esses testes contemplaram:

\begin{itemize}
    \item o login e a identificação do usuário autenticado;
    \item o cadastro e a consulta de aprendentes e de instrumentos avaliativos, incluindo a criação de novas versões e a validação das perguntas;
    \item a restrição de acesso por titularidade em aprendentes, instrumentos, aplicações, relatórios e dashboard;
    \item as regras de finalização das aplicações e a preservação da autoria de aplicações e relatórios;
    \item o controle de acesso por perfil, incluindo a recusa do administrador nas rotas de dados clínicos e na criação e edição de instrumentos;
    \item o registro de auditoria, incluindo a ocultação dos valores clínicos na consulta do administrador;
    \item a transferência de aprendentes entre profissionais e a recusa em excluir ou promover a administrador um profissional que ainda tem aprendentes.
\end{itemize}

Não foram contemplados por testes automatizados o cadastro de usuários, a renovação de sessão, o logout e a recuperação de senha. Como o banco de dados é simulado nesses testes, o comportamento junto ao banco real foi verificado separadamente, por roteiros executados contra uma instância descartável do MySQL, cobrindo o fluxo de aplicação de instrumentos com perguntas de seleção e o registro de auditoria. A separação entre administração e acesso clínico também foi verificada contra o banco real, com usuários e aprendente fictícios criados para esse fim: as rotas clínicas recusaram o administrador, a transferência passou o aprendente ao novo responsável e gerou o registro correspondente em sua auditoria, e nenhuma leitura dos dados da criança foi registrada em nome do administrador.

No frontend (Seção~\ref{sec:frontend}), foram considerados a tela de login, a proteção das rotas conforme o perfil do usuário, o formulário e a tela de visualização dos instrumentos, o preenchimento, o salvamento e a finalização das aplicações, a criação e a impressão dos relatórios, a consulta ao registro de auditoria, o dashboard do administrador e a transferência de aprendentes na página de usuários. Nesses testes, as chamadas à API também são simuladas.

Além dos testes automatizados, foram realizados testes funcionais e exploratórios durante o desenvolvimento, buscando verificar o funcionamento integrado das funcionalidades do sistema. Nessa etapa, foram observados fluxos como autenticação, cadastro e consulta de aprendentes, criação e gerenciamento de instrumentos de avaliação, realização de aplicações, registro de respostas e geração de relatórios. Também foram verificadas situações relacionadas à consistência dos dados e ao acesso às funcionalidades protegidas.

Como complemento aos testes técnicos, foi realizada uma avaliação com profissionais de psicopedagogia, voltada a aspectos que os testes automatizados não alcançam, como dificuldades de compreensão, problemas de navegação e adequação do fluxo de trabalho. Para essa avaliação, foi disponibilizado um questionário estruturado após a utilização do sistema pelos participantes. O instrumento de coleta reuniu informações sobre aspectos como facilidade de uso, organização das informações, clareza das funcionalidades, navegação, adequação das funcionalidades às atividades profissionais e percepção sobre a utilidade do sistema como ferramenta de apoio à avaliação psicopedagógica. Também foi disponibilizado espaço para que os participantes registrassem dificuldades encontradas, sugestões de melhoria e observações sobre as funcionalidades avaliadas.
Os resultados obtidos nessa etapa foram analisados de forma descritiva, considerando as respostas dos participantes e os comentários registrados no questionário. A análise permitiu identificar pontos positivos, limitações e oportunidades de melhoria no sistema. Os resultados quantitativos e qualitativos dessa avaliação são apresentados no Capítulo~\ref{capitulo6}. Essa avaliação constituiu uma etapa de validação da solução desenvolvida, sem caracterizar uma validação clínica ou metodológica dos instrumentos psicopedagógicos utilizados.
